% Generated draft: requires mathematical and source review.
\documentclass[11pt]{article}
\usepackage[margin=25mm]{geometry}
\usepackage{fontspec}
\setmainfont{Latin Modern Roman}
\usepackage{amsmath,amssymb,mathtools,unicode-math}
\setmathfont{Latin Modern Math}
\usepackage{enumitem,needspace,xurl,hyperref}
\setlength{\parindent}{0pt}
\setlength{\parskip}{5pt}
\setlength{\emergencystretch}{3em}
\newcommand{\N}{\mathbb{N}}
\newcommand{\Z}{\mathbb{Z}}
\newcommand{\Q}{\mathbb{Q}}
\newcommand{\R}{\mathbb{R}}
\newcommand{\C}{\mathbb{C}}
\newcommand{\F}{\mathbb{F}}
\newcommand{\A}{\mathbb{A}}
\newcommand{\PP}{\mathbb{P}}
\newcommand{\E}{\mathbb{E}}
\newcommand{\eps}{\varepsilon}
\newcommand{\dd}{\,\mathrm{d}}
\begin{document}
\section*{UnsolvedMath ID: 20002897}
\textbf{Source problem code:} catalogue entry
\subsection{Definitions and mathematical statement}
Fix $k\geq1$. A partition $\lambda$ is $k$-bounded if $\lambda_1\leq k$; write $\ell(\lambda)$ for its number of nonzero parts and $s_\lambda$ for the ordinary Schur function. The workshop asks to prove equivalence of four definitions of $k$-Schur functions, referring to Morse's lecture notes. Those definitions and normalization conventions are not reproduced in this record.

The aid reconstructs the four historical candidates as the original tableau atom $A_\lambda^{(k)}$, the operator/$k$-split candidate $\widetilde A_\lambda^{(k)}$, the inverse weak $k$-Kostka candidate $\widetilde s_\lambda^{(k)}$, and the strong marked-tableau candidate $s_\lambda^{(k)}$. The intended comparison, for all $k$-bounded $\lambda$, is equality of these four families with their appropriate parameter $t$ and common normalization. The weak candidate is defined by an inverse transition matrix, not by directly taking the weak-tableau enumerator. The exact transition basis and matrix are not supplied, so they cannot be silently replaced by a dual family.

According to the supplied literature assessment, a 2020 generic-$t$ unification theorem identifies $\widetilde A_\lambda^{(k)}$, an auxiliary Catalan candidate and the strong $s_\lambda^{(k)}$. At $t=1$ the strong candidate also agrees with the weak inverse. The original-atom comparison and the generic-$t$ comparison with the weak inverse remain open in that assessment. No complete references or formulas for these candidates are included here.

The aid's stated stable-range consequence is precise in its indexing. For nonempty $\lambda$ put
\[
h_M(\lambda)=\lambda_1+\ell(\lambda)-1.
\]
If $h_M(\lambda)\leq k$, then, for generic $t$,
\[
A_\lambda^{(k)}(X;t)=\widetilde A_\lambda^{(k)}(X;t)
=s_\lambda^{(k)}(X;t)=s_\lambda(X).
\]
At $t=1$, $\widetilde s_\lambda^{(k)}$ joins this equality. The empty partition has the usual value $1$. The assertion does not include generic-$t$ equality with the weak inverse outside the supplied results.

The aid's diagnostic example is $k=1$, $\lambda=(1,1)$, $t=1$. In the homology $k$-Schur realization the relevant function is $h_1^2$, whereas the direct affine/dual weak-tableau enumerator is $m_{11}$, with $h_r$ and $m_\mu$ denoting complete and monomial symmetric functions. In characteristic zero,
\[
h_1^2=m_2+2m_{11}.
\]
Even in the usual dual quotient $\Lambda/\langle m_\mu:\mu_1>1\rangle$, its image is $2m_{11}\ne m_{11}$. This is an obstruction to substituting the direct dual enumerator for the inverse weak definition, not a counterexample to the correctly typed four-candidate problem. This partition has main hook $2>k$, so it is outside the preceding stable range.
\subsection{Short English statement}
Compare the four historical $k$-Schur candidates using their original definitions. The supplied aid reports several known comparison edges and equality with ordinary Schur functions in the main-hook stable range, while warning that direct affine/dual enumeration is not the inverse weak candidate.
\subsection{Sources}
\begin{itemize}
\item Local UnsolvedMath registry, permanent ID $20002897$, source code \texttt{catalogue entry}; canonical \texttt{aim-special-functions-notes.json}, \texttt{notes[1]}.
\item AIM workshop \emph{Generalized Kostka polynomials}, Problem (2): \url{https://aimath.org/WWN/kostka/prob2.pdf}. The question refers to Morse's lecture notes. No copy, exact version or bibliographic entry for those notes is supplied; original remarks and literature field are empty.
\item The aid describes the four definitions as historical 2005 candidates and invokes a 2020 generic-$t$ unification theorem, without authors, title or complete authoritative citation for that theorem.
\item Research attempt 1: \texttt{partial\_result}, status \texttt{partially\_solved}, classification \texttt{PARTIAL-PROGRESS}, rigor \texttt{proved}, novelty confidence low. Its comparison graph, hook-stability transfer and minimal example are machine-generated aid, not peer reviewed, and require independent verification before citation.
\end{itemize}
\end{document}
